\documentclass[10pt,a4paper]{amsart}
\usepackage{amsmath,amssymb,amsthm,graphicx}
\usepackage{verbatim}
\newenvironment{reference}{\comment}{\endcomment}
%\newenvironment{reference}{}{}
\newenvironment{slogan}{\comment}{\endcomment}
\newenvironment{history}{\comment}{\endcomment}

% For commutative diagrams we use Xy-pic
\usepackage[all]{xy}

% We use 2cell for 2-commutative diagrams.
\xyoption{2cell}
% Theorem environments.
%
\theoremstyle{plain}
\newtheorem{theorem}[subsection]{Theorem}
\newtheorem{proposition}[subsection]{Proposition}
\newtheorem{lemma}[subsection]{Lemma}

\theoremstyle{definition}
\newtheorem{definition}[subsection]{Definition}
\newtheorem{example}[subsection]{Example}
\newtheorem{exercise}[subsection]{Exercise}
\newtheorem{situation}[subsection]{Situation}

\theoremstyle{remark}
\newtheorem{remark}[subsection]{Remark}
\newtheorem{remarks}[subsection]{Remarks}

\numberwithin{equation}{subsection}

% Macros
%
\def\lim{\mathop{\mathrm{lim}}\nolimits}
\def\colim{\mathop{\mathrm{colim}}\nolimits}
\def\Spec{\mathop{\mathrm{Spec}}}
\def\Hom{\mathop{\mathrm{Hom}}\nolimits}
\def\Ext{\mathop{\mathrm{Ext}}\nolimits}
\def\SheafHom{\mathop{\mathcal{H}\!\mathit{om}}\nolimits}
\def\SheafExt{\mathop{\mathcal{E}\!\mathit{xt}}\nolimits}
\def\Sch{\mathit{Sch}}
\def\Mor{\mathop{\mathrm{Mor}}\nolimits}
\def\Ob{\mathop{\mathrm{Ob}}\nolimits}
\def\Sh{\mathop{\mathit{Sh}}\nolimits}
\def\NL{\mathop{N\!L}\nolimits}
\def\CH{\mathop{\mathrm{CH}}\nolimits}
\def\proetale{{pro\text{-}\acute{e}tale}}
\def\etale{{\acute{e}tale}}
\def\QCoh{\mathit{QCoh}}
\def\Ker{\mathop{\mathrm{Ker}}}
\def\Im{\mathop{\mathrm{Im}}}
\def\Coker{\mathop{\mathrm{Coker}}}
\def\Coim{\mathop{\mathrm{Coim}}}

% Boxtimes
%
\DeclareMathSymbol{\boxtimes}{\mathbin}{AMSa}{"02}

%
% Macros for moduli stacks/spaces
%
\def\QCohstack{\mathcal{QC}\!\mathit{oh}}
\def\Cohstack{\mathcal{C}\!\mathit{oh}}
\def\Spacesstack{\mathcal{S}\!\mathit{paces}}
\def\Quotfunctor{\mathrm{Quot}}
\def\Hilbfunctor{\mathrm{Hilb}}
\def\Curvesstack{\mathcal{C}\!\mathit{urves}}
\def\Polarizedstack{\mathcal{P}\!\mathit{olarized}}
\def\Complexesstack{\mathcal{C}\!\mathit{omplexes}}
% \Pic is the operator that assigns to X its picard group, usage \Pic(X)
% \Picardstack_{X/B} denotes the Picard stack of X over B
% \Picardfunctor_{X/B} denotes the Picard functor of X over B
\def\Pic{\mathop{\mathrm{Pic}}\nolimits}
\def\Picardstack{\mathcal{P}\!\mathit{ic}}
\def\Picardfunctor{\mathrm{Pic}}
\def\Deformationcategory{\mathcal{D}\!\mathit{ef}}
\usepackage{cleveref}
\crefname{originallemma}{original lemma}{original lemmas}
\Crefname{originallemma}{Original lemma}{Original lemmas}
\makeatletter
\newlabel{lemma-stalk-pushforward-closed-immersion}{{04FX}{source}{}{sourcepointer-1}{}}
\newlabel{lemma-stalk-pushforward-closed-immersion@cref}{{[originallemma][0][]04FX}{[1][0][]source}{}{}{}}
\makeatother
\crefname{originaltheorem}{original theorem}{original theorems}
\Crefname{originaltheorem}{Original theorem}{Original theorems}
\makeatletter
\newlabel{theorem-exactness-stalks}{{03PU}{source}{}{sourcepointer-2}{}}
\newlabel{theorem-exactness-stalks@cref}{{[originaltheorem][0][]03PU}{[1][0][]source}{}{}{}}
\makeatother
\crefname{originallemma}{original lemma}{original lemmas}
\Crefname{originallemma}{Original lemma}{Original lemmas}
\makeatletter
\newlabel{sites-lemma-exactness-properties}{{04D6}{source}{}{sourcepointer-3}{}}
\newlabel{sites-lemma-exactness-properties@cref}{{[originallemma][0][]04D6}{[1][0][]source}{}{}{}}
\makeatother
\crefname{originallemma}{original lemma}{original lemmas}
\Crefname{originallemma}{Original lemma}{Original lemmas}
\makeatletter
\newlabel{lemma-stalk-pullback}{{03Q1}{source}{}{sourcepointer-4}{}}
\newlabel{lemma-stalk-pullback@cref}{{[originallemma][0][]03Q1}{[1][0][]source}{}{}{}}
\makeatother
\title{The pushforward along a closed immersion is fully faithful, with essential image described by support}
\author{The Stacks Project Authors}
\date{stacks-project:etale-cohomology.tex:local-snapshot-20261011; GNU FDL 1.2 or later}
\begin{document}
\maketitle
\noindent\textit{Editorial note.} All delivered mathematical text is the author's own, copied from the pinned Stacks Project chapter file etale-cohomology.tex; no statement or proof step was rewritten, reordered or removed.\par\medskip
\noindent\textit{Editorial note.} The proof is the author's complete argument for both parts of the statement and for the final assertion about the inverse image.\par\medskip
\noindent\textit{Editorial note.} Every labelled reference occurring inside the counted spans is carried in assembly.external\_references as a non-load-bearing navigation pointer (same-chapter ones at their definition site, cross-chapter ones at the invocation inside the counted span), so no load-bearing mathematics is replaced by a stub. No bibliography citation occurs inside the counted spans.\par\medskip
\noindent\textit{Editorial note.} The checked chapter list for same-name collisions is the set of stacks chapters staged in this repository: chow, cohomology, divisors, etale-cohomology, flat, more-algebra, more-morphisms, morphisms, spaces-more-morphisms, varieties together with this card's chapter etale-cohomology.\par\medskip
\noindent\textit{Editorial note.} The preamble is the chapter's own native preamble excerpt already used for this source: the Xy-pic setup (native preamble lines 14-18) and the native theorem declarations with the mathematical macros (native preamble lines 154-217).\par\medskip
\noindent\textit{Allowed standard prerequisites (editorial).}\begin{itemize}
\item Closed immersion and the comparison of the etale sites of X and of the complement. Same chapter, the open immersion of the complement of Z and the resulting picture in which a sheaf on X\_etale is determined by its restriction to Z and to the complement. Supplies the two functors compared and the description of the essential image.
\item Adjunction between the small inverse image and pushforward along a closed immersion. Same chapter, the adjunction i\_\{small\}\textasciicircum{}\{-1\} -| i\_\{small,*\} and the fact that the inverse image is a left inverse of the pushforward. Supplies the fully faithfulness and the last assertion of the statement.
\end{itemize}\medskip
\ifdefined\hypertarget\hypertarget{sourcepointer-1}{}\fi
\noindent\textit{Source reference.} Original lemma 04FX: sources/stacks-etale-cohomology/pinned-original-source.tex, source lines 6644--6644. Navigation pointer invoked at proof line 6734; the counted span contains the pointer only and the cited result is not reproduced inside it, so the pointer is registered as a non-load-bearing navigation stub in the established form for Stacks material. No bibliography citation is involved.\par\medskip
\ifdefined\hypertarget\hypertarget{sourcepointer-2}{}\fi
\noindent\textit{Source reference.} Original theorem 03PU: sources/stacks-etale-cohomology/pinned-original-source.tex, source lines 3546--3546. Navigation pointer invoked at proof line 6736; the counted span contains the pointer only and the cited result is not reproduced inside it, so the pointer is registered as a non-load-bearing navigation stub in the established form for Stacks material. No bibliography citation is involved.\par\medskip
\ifdefined\hypertarget\hypertarget{sourcepointer-3}{}\fi
\noindent\textit{Source reference.} Original lemma 04D6: sources/stacks-etale-cohomology/pinned-original-source.tex, source lines 6738--6738. Navigation pointer invoked at proof line 6738; the counted span contains the pointer only and the cited result is not reproduced inside it, so the pointer is registered as a non-load-bearing navigation stub in the established form for Stacks material. No bibliography citation is involved.\par\medskip
\ifdefined\hypertarget\hypertarget{sourcepointer-4}{}\fi
\noindent\textit{Source reference.} Original lemma 03Q1: sources/stacks-etale-cohomology/pinned-original-source.tex, source lines 4649--4649. Navigation pointer invoked at proof line 6748; the counted span contains the pointer only and the cited result is not reproduced inside it, so the pointer is registered as a non-load-bearing navigation stub in the established form for Stacks material. No bibliography citation is involved.\par\medskip
\begin{proposition}
\label{proposition-closed-immersion-pushforward}
Let $i : Z \to X$ be a closed immersion of schemes.
\begin{enumerate}
\item The functor
$$
i_{small, *} :
\Sh(Z_\etale)
\longrightarrow
\Sh(X_\etale)
$$
is fully faithful and its essential image is those sheaves of sets
$\mathcal{F}$ on $X_\etale$ whose restriction to $X \setminus Z$ is
isomorphic to $*$, and
\item the functor
$$
i_{small, *} :
\textit{Ab}(Z_\etale)
\longrightarrow
\textit{Ab}(X_\etale)
$$
is fully faithful and its essential image is those abelian sheaves on
$X_\etale$ whose support is contained in $Z$.
\end{enumerate}
In both cases $i_{small}^{-1}$ is a left inverse to the functor
$i_{small, *}$.
\end{proposition}

\begin{proof}
Let's discuss the case of sheaves of sets.
For any sheaf $\mathcal{G}$ on $Z$ the morphism
$i_{small}^{-1}i_{small, *}\mathcal{G} \to \mathcal{G}$
is an isomorphism by
Lemma \ref{lemma-stalk-pushforward-closed-immersion}
(and
Theorem \ref{theorem-exactness-stalks}).
This implies formally that $i_{small, *}$ is fully faithful, see
Sites, Lemma \ref{sites-lemma-exactness-properties}.
It is clear that $i_{small, *}\mathcal{G}|_{U_\etale} \cong *$
where $U = X \setminus Z$. Conversely, suppose that $\mathcal{F}$
is a sheaf of sets on $X$ such that $\mathcal{F}|_{U_\etale} \cong *$.
Consider the adjunction mapping
$$
\mathcal{F} \longrightarrow i_{small, *}i_{small}^{-1}\mathcal{F}
$$
Combining
Lemmas \ref{lemma-stalk-pushforward-closed-immersion} and
\ref{lemma-stalk-pullback}
we see that it is an isomorphism. This finishes the proof of (1).
The proof of (2) is identical.
\end{proof}

\end{document}
